\documentclass{article}
\usepackage{graphicx}
\usepackage{amsmath}
\usepackage{amsfonts}
\usepackage{bookmark}
\usepackage{booktabs}
\usepackage{exercises}
\graphicspath{ {} }
\usepackage{indentfirst}
\usepackage{amsthm}
\usepackage{yhmath}
\usepackage{geometry}
\usepackage{placeins}
\usepackage{csquotes}
\usepackage{pgfplots}
\newcommand*{\oldneg}{\mathord{\sim}}
\newcommand\showdiv[1]{\overline{\smash{)}#1}}
\newcommand\Mydiv[1]{%
\smash{\raise.3ex\hbox{$\big)$}}\mkern-8mu
        \overline{\enspace\strut#1}}
\setlength{\tabcolsep}{5pt}
\theoremstyle{definition}
\newtheorem{definition}{Definition}

\setlength{\parindent}{.5in}
\setlength{\parskip}{12pt}
\renewcommand{\arraystretch}{2}
\geometry{margin=1in}
\newcommand{\?}{\stackrel{?}{=}}
\pgfplotsset{width=7cm,compat=1.18}\usepgfplotslibrary{statistics}

\title{6.3 Answers}
\begin{document}
\maketitle

\section*{1}
A normal distribution is a bell-shaped curve that is symmetric about the mean.

\section*{2}
To find $P(\le 1.4)$, we first go down to the row for $1$, and then finding where that row intersects the column of $.4$. This point of intersection is $P(\le 1.4)$.

\section*{7}
$P(\bar{x}-\sigma\le x\le\bar{x}+\sigma)=34+34=68$.

\section*{8}
$P(\bar{x}-3\sigma\le x\le\bar{x})=2.35+13.5+34=49.85$

\section*{9}
$13.5+2.35+0.15=16$

\section*{10}
$2.35+2.35=4.7$

\section*{11}
The probabilities between 29 and 37 are those in the range $\bar{x}-\sigma\le x\le \bar{x}+\sigma$, meaning the probability is $34+34=68$.

\section*{15}
The range here is $\bar{x}-\sigma\le x$ meaning the probability is $34+34+13.5+2.35+0.15=84$.

\section*{16}
The range here is $\bar{x}+\sigma> x$ meaning the probability is $0.15+2.35+13.5+34=50$.

\section*{17}
C

\section*{18}
B

\section*{20}
We find the $z$ score for when $x=80$: \[ z=\frac{80-64}{7}\approx 2.3 \] which means $P(z\le 2.3)$ corresponds to a probability of $.9893$ in the standard normal table.

\section*{26}
We find the $z$ score for when $x=60$: \[ z=\frac{60-64}{7}\approx -0.6 \] and the $z$ score for when $x=75$ is \[ z=\frac{75-74}{7}\approx 0.1 \] Thus we are finding the probability for
$P(-0.6\le z\le 0.1)$, which, from the standard normal table, is between $.2743$ to $.5398$.

\section*{28}
We are considering the 84th percentile for normally distributed data, that is, a number such that $84\%$ of the data are less than or equal to that number. The number that represents this value is 
$\bar{x}+\sigma$.

\section*{29}
The error here is that the computed probability is for values less than or equal to $3.8$ bushels, not at least. Instead, to find the probability for being at least, we subtract the probability for being less than or 
equal to $3.8$ from 1: $1-0.2119=0.7881$. Hence the probability that a plot yielded at least 3.8 bushels is $0.7881$.

\section*{33}
(a): The $z$-score for 19.4 is $\frac{19.4-20}{0.25}=-2.4$ while the $z$-score for 20.4 is $\frac{20.4-20}{0.25}=1.6$.


(b): This can be found by finding the probability for $z=-2.4$ ($x=19.4) in the standard normal table, which is $.0082$.


(c): First, the probability that a box weight at most 20.4 ($z=1.6$) is, by the standard normal table, $.9452$. So, to find the probability that a box is between 19.4 and 20.4 ounces, we subtract the probability that it is 
less than or equal to 19.4 ounces from the probability that it is less than or equal to 20.4 ounces. Hence the probability is $.9452-.0082=.937$. 

\section*{35}
(a): Lisa's test score has a $z$-score of $\frac{30-20}{4.2}\approx 2.4$


(b): Ann's test score has a $z$-socre of $\frac{610-500}{90}\approx 1.2$


(c): The student who scored better on the score would be the one who got a score that people are more likely to score less than or equal to. This can be found by the standard normal table. The probability for 
Lisa is about $.9918$ while the probability for Ann is $.8849$, meaning Lisa got the better score.

\section*{36}
(a): We first find the $z$-score for being 6-feet (72 inches), which is $\frac{72-69}{2.75}\approx 1.1$ meaning the probability for being less than or equal to 6 feet is 
about $.8643$. Hence the probability that you are more than 6 feet tall is $1-0.8643=0.1357$. 
Then we set up a binomial experiment where the chance of success is $.1357$. Hence the probability that, from an experiment of $3$ men, all three are more than 6 feet, is:
\[ .1357^3\approx .002 \] or about $0.2\%$.


(b): This is a similar method as to the first part. The $z$-score for being $65$ inches is $\frac{65-69}{2.75}\approx -1.5$ while the $z$-score for being $75$ inches is $\frac{75-69}{2.75}\approx 2.2$. Hence the 
probabilities for being less than or equal to 65 inches and 75 inches are $.668$ and $.9821$, respectively. Hence the probability that someone has a height between 75 and 65 inches is 
$.9821-.668=0.3141$. Thus, we set up a binomial experiment where the chance of success is $0.3141$. Hence the probability that, from an experiment of $5$ men, all five are between 65 and 75 inches is:
\[ .3141^5\approx .003 \] or about $0.3\%$. 

\end{document}
